\chapter{자료와 재현 근거}

\section{AMD/Xilinx 문서}

\begin{itemize}
  \item \href{https://docs.amd.com/r/en-US/ug901-vivado-synthesis}{UG901, Vivado Design Suite User Guide: Synthesis}: RAM/DSP/SRL coding과 synthesis option
  \item \href{https://docs.amd.com/r/en-US/ug949-vivado-design-methodology}{UG949, UltraFast Design Methodology Guide}: timing path 분석, pipelining, fanout, control set, inference
  \item \href{https://docs.amd.com/r/en-US/ug906-vivado-design-analysis}{UG906, Vivado Design Suite User Guide: Design Analysis and Closure Techniques}: timing/utilization report 해석
  \item \href{https://docs.amd.com/v/u/en-US/ug479_7Series_DSP48E1}{UG479, 7 Series DSP48E1 Slice User Guide}: DSP48E1 data path와 internal register
  \item \href{https://docs.amd.com/v/u/en-US/ug473_7Series_Memory_Resources}{UG473, 7 Series FPGAs Memory Resources}: RAMB36E1/RAMB18E1 구성
  \item \href{https://docs.amd.com/r/en-US/ug474_7Series_CLB}{UG474, 7 Series FPGAs Configurable Logic Block User Guide}: LUT, FF, carry, slice 구조
\end{itemize}

\section{사례 RTL 파일}

\begin{longtable}{L{58mm}L{104mm}}
\toprule
파일 & 이 책에서 확인한 내용 \\
\midrule
\rtlfile{MDL\_XXX\_ntt\_core.v} & BRAM command, physical tag, occupancy, butterfly/writeback alignment \\
\rtlfile{MDL\_XXX\_layout\_mapper.v} & side/compact-row bijection과 고정 affine tap \\
\rtlfile{MDL\_XXX\_bram\_dp.v} & true-dual-port synchronous BRAM inference template \\
\rtlfile{MDL\_XXX\_zeta\_rom.v} & packed 2048$\times$32 two-read-port twiddle memory \\
\rtlfile{MDL\_XXX\_XXX\_unifiedMUL.v} & two-DSP, six-stage ML-DSA/R16 shared multiplier \\
\rtlfile{MDL\_XXX\_unifiedBUT.v} & two multiplier instance와 unified butterfly pipeline \\
\rtlfile{MDL\_XXX\_ntt\_addrgen.v} & logical address schedule와 NTRU+864 special mode \\
\rtlfile{MDL\_XXX\_ntt\_engine.v} & transport-independent core boundary \\
\rtlfile{MDL\_XXX\_axis\_load.v} & 64-bit beat의 2/4-coefficient write \\
\rtlfile{MDL\_XXX\_axis\_dump.v} & finite FIFO credit와 AXI backpressure \\
\bottomrule
\end{longtable}

\section{측정 report}

수치는 companion project에 보존된 다음 report를 source of truth로 삼았다.

\begin{itemize}
  \item 최초 100~MHz OOC: \rtlfile{reports/vivado\_impl\_baseline}
  \item 200~MHz 목표판: \rtlfile{reports/vivado\_200mhz}와
        \rtlfile{doc/200mhz\_result.md}
  \item 최종 세 범위: \rtlfile{reports/memory\_conflict\_mc/axis},
        \rtlfile{bram}, \rtlfile{mdl\_only}
  \item 최종 board 측정:
        \rtlfile{board\_benchmark\_150mhz\_final\_20260903.md/csv}
\end{itemize}

표를 다른 문서로 옮길 때는 report path, Vivado version, device, clock constraint,
top-level을 함께 옮긴다. 숫자만 복사하면 이후의 설계 변경과 비교 범위를 추적할
수 없다.

\section{개정 원칙}

이 경험서는 FPGA 기본 자원과 Verilog-to-hardware mapping을 먼저 독립적으로
설명한다. 각 설명은 ``RTL 표현--생성되는 회로--면적 근사--timing 경로--검증
조건''을 연결한다. Unified NTT의 RTL과 report는 일반 원칙을 실제 개발에서
확인한 사례로 사용한다. 프로젝트 고유 신호 이름과 결과표만으로 본문을 전개하지
않으며, 독자가 다른 RTL에도 같은 판단 순서를 적용할 수 있게 한다.

구조 기반 근사와 실측값은 명확히 구분한다. 연산자 micro-benchmark는 동일한
device와 synthesis 조건에서 자원 mapping을 비교하는 용도이고, NTT post-route
path는 logic/route/carry가 결합된 실제 timing 경험을 보여 주는 용도다. 특정
LUT나 adder에 보편적인 고정 ns 값을 부여하지 않는다.
